\section{Structural Identities and Obstruction Tools}
\label{sec:structural-tools}

The preceding lower bounds use several small pieces of structure that are
useful beyond any one hard family.  This section collects them in one place.
The first two tools turn symmetry and cross-instance averaging into invariant
certificates.  The third says how much matrix rank an SDP Newton step must
carry to make complementarity progress.  The remaining tools delimit what
sparse support, bounded treewidth, and geometric locality can and cannot buy.
They are supporting results: most of their ingredients are classical, and we
separate the new obstruction formulations from those ingredients throughout.

\subsection{Symmetry as virtual work}
\label{sec:structural-symmetry}

Work in a finite-dimensional Euclidean space.  Let $R$ be an orthogonal
involution, let $\Omega$ be an open convex set with $R\Omega=\Omega$, and let
$F:\Omega\to\R$ be differentiable, convex, and $R$-invariant.  Assume also
that $R^Tc=c$ and use the stationarity convention
\begin{equation}
 A^Ty=c+\mu\nabla F(x),\qquad \mu>0.                 \label{eq:structural-stationarity}
\end{equation}
Thus the equality term in the Lagrangian is $-y^T(Ax-b)$.  Put
\begin{equation}
 x_-:=\frac{I-R}{2}x,\qquad d:=x-Rx=2x_-,\qquad r_-:=Ax_-.
                                                               \label{eq:structural-antisymmetric-source}
\end{equation}
The definition $r_-=Ax_-$, rather than a selected block of $b$, is what makes
the following statement insensitive to mixed equality rows.

\begin{theorem}[Involution work identity]
\label{thm:structural-involution-work}
Under the preceding assumptions,
\begin{align}
 r_-^Ty
 &=\mu\inner{x_-}{\nabla F(x)}                                      \notag\\
 &=\frac{\mu}{4}\inner{x-Rx}{\nabla F(x)-\nabla F(Rx)}             \label{eq:structural-work-chain}\\
 &=\frac{\mu}{2}D_F(Rx,x)
  =\frac{\mu}{2}D_F(x,Rx)\geq0,                                    \notag
\end{align}
where $D_F(u,v)=F(u)-F(v)-\inner{\nabla F(v)}{u-v}$.
If $F\in C^2(\Omega)$ and
\begin{equation}
 \overline H_x:=\int_0^1\nabla^2F(Rx+td)\,dt,
\end{equation}
then, exactly,
\begin{equation}
 r_-^Ty=\frac{\mu}{4}d^T\overline H_xd
       =\mu x_-^T\overline H_xx_-.                         \label{eq:structural-integrated-hessian}
\end{equation}
For a standard self-concordant barrier, set
$x_+=(x+Rx)/2$ and $\delta=\norm{x_-}_{x_+}<1$.  Then the
integrated chord comparison gives sharper local constants than a worst-point
comparison, and the universal lower bound can be strengthened further:
\begin{equation}
 \boxed{
 \mu\left(1-\delta+\frac{\delta^2}{3}\right)\delta^2
 \leq\frac{\mu\delta^2}{1+\delta}
 \leq r_-^Ty\leq\frac{\mu\delta^2}{1-\delta}.}
                                                               \label{eq:structural-sharp-chord}
\end{equation}
\end{theorem}

\begin{proof}
The invariant objective does no antisymmetric work:
\[
 \inner{x_-}{c}=\tfrac12\inner{x}{c-R^Tc}=0.
\]
Pairing \eqref{eq:structural-stationarity} with $x_-$ gives the first
equality in \eqref{eq:structural-work-chain}.  Differentiating
$F(Rx)=F(x)$ gives $\nabla F(Rx)=R\nabla F(x)$.  Since $Rd=-d$,
\[
 \inner{d}{\nabla F(Rx)}
 =\inner{Rd}{\nabla F(x)}=-\inner{d}{\nabla F(x)},
\]
which proves the second equality.  Gradient monotonicity proves
nonnegativity.  Invariance also gives
\[
 D_F(Rx,x)=D_F(x,Rx)=\inner{\nabla F(x)}{x-Rx},
\]
so the Bregman equalities follow.  The fundamental theorem of calculus on
the chord gives \eqref{eq:structural-integrated-hessian}.

For the last claim, parameterize the chord as $x_++ux_-$,
$-1\leq u\leq1$.  Self-concordant Hessian comparison bounds its directional
energy by
\[
 (1-|u|\delta)^2\delta^2
 \leq x_-^T\nabla^2F(x_++ux_-)x_-
 \leq(1-|u|\delta)^{-2}\delta^2.
\]
The relevant chord averages are
\[
 \int_0^1(1-u\delta)^2\,du
 =1-\delta+\delta^2/3,
 \qquad
 \int_0^1(1-u\delta)^{-2}\,du=(1-\delta)^{-1},
\]
so these constants are the exact integrals of the comparison envelopes, not
generally attained bounds.  For the stronger lower bound, put
\[
 \phi(u)=x_-^T\nabla^2F(x_++ux_-)x_-.
\]
Self-concordance gives
$|(\phi^{-1/2})'|\leq1$ and $\phi(0)=\delta^2$, hence
$\phi(u)\geq\delta^2/(1+u\delta)^2$ for $0\leq u\leq1$.
Invariance makes $\phi$ even and
$r_-^Ty=\mu\int_0^1\phi(u)\,du$, which proves the middle lower bound.
It strictly dominates the chord-comparison lower constant for
$0<\delta<1$, since
\[
 \frac1{1+\delta}-\left(1-\delta+\frac{\delta^2}{3}\right)
 =\frac{\delta^2(2-\delta)}{3(1+\delta)}>0.
\]
\end{proof}

No rank assumption on $A$ is hidden here.  If $\widetilde y$ is another
multiplier, then $A^T(\widetilde y-y)=0$ and hence
$(Ax_-)^T(\widetilde y-y)=0$.  Moreover, under any invertible row change
\begin{equation}
 A'=UA,\qquad y'=U^{-T}y,\qquad r_-'=Ur_-,
 \qquad (r_-')^Ty'=r_-^Ty.                              \label{eq:structural-row-covariance}
\end{equation}
Only the source--multiplier pairing is invariant; an ill-conditioned $U$
can shrink a multiplier norm while enlarging the source or the input
normalization.  If rows genuinely split as $A_+R=A_+$ and $A_-R=-A_-$
and $Ax=b=(b_+,b_-)$, then $r_-=(0,b_-)$ and the pairing is the familiar
$b_-^Ty_-$.  Arbitrary
row mixing preserves \eqref{eq:structural-row-covariance}, but not that
coordinate shortcut.

Two specializations fix the constants used earlier.  For paired positive
variables $x=(u,v)$, $R(u,v)=(v,u)$, and
$F=-\sum_j(\log u_j+\log v_j)$, put $d=u-v$ and $q=u+v$.  Then
\begin{equation}
 r_-^Ty=\frac{\mu}{2}\sum_j\frac{d_j^2}{u_jv_j}
       =\sum_j\frac{2\mu d_j^2}{q_j^2-d_j^2}.             \label{eq:structural-paired-log}
\end{equation}
This is exactly \eqref{eq:frontier-central-work}; the proof in
\cref{thm:frontier-work} is the paired-coordinate specialization, not a
second general identity.

For the log-determinant barrier on $\Sym_{++}^r$, let
$\mathcal R(X)=QXQ$ for a symmetric orthogonal $Q$, and set
$X_-=(X-QXQ)/2$.  With the Frobenius pairing,
\begin{align}
 (\mathcal A X_-)^Ty
 &=\frac{\mu}{2}\bigl[\tr(QXQX^{-1})-r\bigr]              \label{eq:structural-stein}\\
 &=\mu\int_0^1
 \norm{Z_t^{-1/2}X_-Z_t^{-1/2}}_F^2\,dt,
 \quad Z_t=QXQ+t(X-QXQ).                                  \label{eq:structural-logdet-integral}
\end{align}
The bracketed trace quantity in the first line is half the symmetrized Stein
loss, so the work is $\mu/4$ times that symmetrization.  This is not the
distinct Jensen--Bregman LogDet, or S-divergence.  The second line is the
exact noncommutative integrated-Hessian formula; no commuting assumption is
used.  If $\ell I\preceq X\preceq UI$, it lies between
$\mu\norm{X_-}_F^2/U^2$ and $\mu\norm{X_-}_F^2/\ell^2$.
Here $U$ is a spectral endpoint, unrelated to the row-change matrix in
\eqref{eq:structural-row-covariance}.

At a fixed start $x^0=Rx^0$, the Newton version is equally direct.  If
$G_0\succ0$ commutes with $R$ and
\begin{equation}
 A\Delta x=r_p,\qquad G_0\Delta x-A^T\Delta y=h_+,
 \qquad Rh_+=h_+,
\end{equation}
then, for $\Delta x_-=(I-R)\Delta x/2$,
\begin{equation}
 \boxed{(A\Delta x_-)^T\Delta y
 =\Delta x_-^TG_0\Delta x_-\geq0.}                       \label{eq:structural-newton-work}
\end{equation}
This recovers \eqref{eq:frontier-newton-work}.  A nonzero antisymmetric
forcing $h_-$ adds $-\inner{\Delta x_-}{h_-}$ and destroys the sign.

The involution is only the two-element case.  Let a finite or compact group
$\mathcal G$ act orthogonally, and assume $g\Omega=\Omega$ for every
$g\in\mathcal G$.  Let
\begin{equation}
 P_0=\int_{\mathcal G}g\,d\nu(g),\qquad
 x_0=P_0x,\qquad x_\perp=(I-P_0)x                       \label{eq:structural-reynolds}
\end{equation}
be the Reynolds projection.  The orbit is compact, and Carath\'eodory's
theorem shows that its convex hull is a compact subset of the convex domain;
thus $x_0\in\Omega$ and the segment from $x_0$ to $x$ stays in $\Omega$.

\begin{theorem}[Compact-group work]
\label{thm:structural-group-work}
If $F$ and $c$ are $\mathcal G$-invariant and
\eqref{eq:structural-stationarity} holds, then
\begin{equation}
 (Ax_\perp)^Ty
 =\mu\inner{x_\perp}{\nabla F(x)-\nabla F(x_0)}\geq0.
                                                               \label{eq:structural-group-central}
\end{equation}
When the Hessian exists, the complete chain is
\begin{align}
 (Ax_\perp)^Ty
 &=\mu\inner{x_\perp}{\nabla F(x)-\nabla F(x_0)}          \notag\\
 &=\mu\int_0^1\inner{x_\perp}
 {\nabla^2F(x_0+tx_\perp)x_\perp}\,dt                     \notag\\
 &\geq0.                                                  \label{eq:structural-group-hessian}
\end{align}
At a group-fixed Newton start, suppose $G_0\succeq0$ commutes with the action and
\begin{equation}
 A\Delta x=r_p,\qquad
 G_0\Delta x-A^T\Delta y=h_0,\qquad gh_0=h_0
 \quad(g\in\mathcal G).                                   \label{eq:structural-group-newton-system}
\end{equation}
Define $\Delta x_\perp=(I-P_0)\Delta x$.  Then
\begin{equation}
 (A\Delta x_\perp)^T\Delta y
 =\Delta x_\perp^TG_0\Delta x_\perp\geq0.               \label{eq:structural-group-newton}
\end{equation}
\end{theorem}

\begin{proof}
Both $c$ and $\nabla F(x_0)$ are fixed by the group and hence orthogonal to
$x_\perp$.  Pair stationarity with $x_\perp$, then apply gradient
monotonicity between $x_0$ and $x$.  Integration gives
\eqref{eq:structural-group-hessian}.  The Newton identity follows by pairing
the stationarity equation with the nontrivial-representation component;
$G_0$ preserves the orthogonal representation decomposition.
\end{proof}

For coordinate permutations, $P_0$ replaces each orbit by its average; for
the logarithmic barrier the work on an orbit $O$ is
$\mu(\bar x_O\sum_{i\in O}x_i^{-1}-|O|)$.  The simplex gives the same
formula with $\bar x=1/n$.  For the standard $SO(k)$ action and the unit-ball
barrier, it is $2\mu\norm z^2/(1-\norm z^2)$.  These examples explain the
permutation and rotation gadgets excluded in \cref{sec:frontier-synthesis}.

\begin{remark}[Scope and provenance]
\label{rem:structural-work-scope}
The sign reverses with the opposite multiplier convention.  An asymmetric
objective or Newton forcing contributes an uncontrolled linear term.  Mixed
rows require the canonical source $Ax_-$; convexity without strict curvature
gives only nonnegativity.  The scalar identity extends to a nonorthogonal
linear involution by using $\nabla F(Rx)=R^{-T}\nabla F(x)$, but the Euclidean
orthogonal decomposition does not.  Finally, arbitrary edge transports need
not arise from one simultaneous public symmetry action and may have no public
Reynolds projector.  Even when they do come from a common representation,
the Reynolds identity supplies no cut-local conservation law; the
signed/gain incidence and holonomy analysis of \cref{sec:sdp-hardness} is
separate.

The divergence in \eqref{eq:structural-work-chain} is the classical
symmetrized Bregman divergence
\cite{structural-bregman1967,structural-nielsen2009-bregman}; Reynolds
reduction and symmetric criticality are likewise classical
\cite{structural-palais1979-symmetric-criticality}.  The contribution used
in this paper is narrower: the row-mixing-invariant source
\eqref{eq:structural-antisymmetric-source} is an obstruction template for
sparse QIPM gadgets.
\end{remark}

\subsection{Central mixtures in semidefinite programming}
\label{sec:structural-sdp-mixtures}

The scalar Kantorovich argument in \cref{sec:frontier-mixture-supplements}
has a fully noncommutative analogue.  Consider SDPs with common $C,b$ and
input-dependent measurement maps $\mathcal A_\xi$, with matrix order $r\geq1$
and $N\geq1$ neighboring inputs.  Fix one Frobenius-isometric
$\operatorname{svec}:\Sym^r\to\R^{r(r+1)/2}$.  All sparsity and coefficient
bounds below refer to the resulting matrix $A_\xi$, not to the unvectorized
data.  At a base input $\varsigma$, suppose every one-bit neighbor
$\varsigma^{(i)}$ has an exact center at the same $\mu>0$:
\begin{equation}
 \mathcal A_{\varsigma^{(i)}}(X_i)=b,\qquad
 \mathcal A_{\varsigma^{(i)}}^*(y_i)+S_i=C,\qquad
 S_i=\mu X_i^{-1},\qquad X_i\succ0.                       \label{eq:structural-sdp-centers}
\end{equation}
For weights $w_i\geq0$, $\sum_iw_i=1$, write
\begin{equation}
 \begin{aligned}
 X_w&=\sum_iw_iX_i,& y_w&=\sum_iw_iy_i,&
 S_w&=\sum_iw_iS_i,\\
 D_i&=X_i-X_w,&&&
 Z_w&=X_w^{1/2}(S_w/\mu)X_w^{1/2}.
 \end{aligned}                                            \label{eq:structural-sdp-average}
\end{equation}

\begin{lemma}[Matrix multiplicative variance]
\label{lem:structural-matrix-variance}
The following identity is exact:
\begin{equation}
 \boxed{
 Z_w-I=X_w^{-1/2}\left(\sum_iw_iD_iX_i^{-1}D_i\right)X_w^{-1/2}.}
                                                               \label{eq:structural-matrix-variance}
\end{equation}
Consequently $Z_w\succeq I$, with equality exactly when all positively
weighted $X_i$ equal $X_w$.
\end{lemma}

\begin{proof}
Expanding the middle term and using $\sum_iw_iX_i=X_w$ gives
\[
 \sum_iw_iD_iX_i^{-1}D_i
 =X_w\left(\sum_iw_iX_i^{-1}\right)X_w-X_w.
\]
Congruence by $X_w^{-1/2}$ proves the identity.  Every un-congrued summand is
positive semidefinite, and a positive weighted sum of them vanishes only
when each corresponding $D_i$ vanishes.
\end{proof}

Identity \eqref{eq:structural-matrix-variance} is an elementary exact
expansion of the arithmetic--harmonic defect; it is not claimed as a new
matrix inequality.

\begin{theorem}[Noncommutative central sandwich]
\label{thm:structural-sdp-kantorovich}
If $0<m\leq M$, $mI\preceq X_i\preceq MI$ for every $i$, and $R=M/m$,
then
\begin{equation}
 \boxed{I\preceq Z_w\preceq K(R)I,\qquad
 K(R)=\frac{(R+1)^2}{4R}.}                                \label{eq:structural-sdp-sandwich}
\end{equation}
The constant is sharp, already for an equal scalar mixture of $m$ and $M$.
No two matrices in the mixture need commute.
\end{theorem}

Here and below $R=M/m$ is a spectral ratio, unrelated to the involution $R$
in \cref{sec:structural-symmetry}.

\begin{proof}
The lower bound is \cref{lem:structural-matrix-variance}.  Functional
calculus, applied to one $X_i$ at a time, gives
\begin{equation}
 X_i+mM X_i^{-1}\preceq(m+M)I.
\end{equation}
Averaging and congruencing by $X_w^{1/2}$ yields
\[
 Z_w\preceq\frac{(m+M)X_w-X_w^2}{mM}.
\]
Every eigenvalue $\lambda$ of $X_w$ is in $[m,M]$, and
\[
 \frac{\lambda(m+M-\lambda)}{mM}
 \leq\frac{(m+M)^2}{4mM}=K(R).
\]
Only $X_i$ and $X_w$ individually entered functional calculus, so no
commutation was assumed.
\end{proof}

The variance identity supplies a dimension-free refinement of the spectral
ratio bound.  If only $X_i\succeq mI$ with $m>0$ is assumed, then
\begin{align}
 \norm{Z_w-I}_F
 &\leq\frac1{m^2}\sum_iw_i\norm{D_i}_F^2,                 \label{eq:structural-sdp-variance-frobenius}\\
 \norm{Z_w-I}_{\mathrm{op}}
 &\leq\frac1{m^2}\sum_iw_i\norm{D_i}_{\mathrm{op}}^2.   \label{eq:structural-sdp-variance-operator}
\end{align}
Indeed, the two inverse factors come separately from the outer congruence
and $X_i^{-1}$ in \eqref{eq:structural-matrix-variance}.  Multiplying by
$\mu$ gives the corresponding complementarity-residual bounds.  They survive
standard recentering: if
\begin{equation}
 \mu_w=\frac{\inner{X_w}{S_w}}r=\mu\alpha,\qquad
 \alpha=\frac{\tr Z_w}{r}\geq1,
\end{equation}
then
\[
 \mu_w^{-1}X_w^{1/2}S_wX_w^{1/2}-I
 =\frac{(Z_w-I)-[\tr(Z_w-I)/r]I}{\alpha}.
\]
Removing the trace component is an orthogonal projection in Frobenius norm,
and it does not enlarge the operator range because $Z_w-I\succeq0$.

The objective identity uses no feasibility of the average.  Every point in
\eqref{eq:structural-sdp-centers} has algebraic gap $r\mu$, so linearity gives
\begin{equation}
 \boxed{\inner C{X_w}-b^Ty_w=r\mu.}                       \label{eq:structural-sdp-gap}
\end{equation}
By contrast,
$\inner{X_w}{S_w}-[\inner C{X_w}-b^Ty_w]
=\mu\tr(Z_w-I)$ is nonnegative.  Under the common sandwich
$0<m\leq M$ and $mI\preceq X_i\preceq MI$, it is at most
$r[K(M/m)-1]\mu$.  This is an arithmetic--harmonic Jensen gap at a
generally base-infeasible point, not an application of weak duality.

We next combine the operator statement with the sparse single-flip residual.
Let $s_r,s_c$ bound the row and column sizes of one fixed union support
of the matrices $A_\xi$, let
$M_{\rm dep}$ count input-dependent positions, each of which is a function
of one designated bit $\varsigma_{\ell(r,c)}$ (a bit may designate several
positions), let their magnitudes be at most $B\geq0$, and assume that every
coordinate of
$\operatorname{svec}(X_i),y_i,\operatorname{svec}(S_i)$ has magnitude at
most $H\geq0$.  Set
\begin{equation}
 s_{\rm KKT}=\max\{s_r,s_c+1\},\qquad
 B_{\rm KKT}=\max\{B,1\}.
\end{equation}
For the uniform average of all $N$ one-bit neighbors, the same incidence
count as \cref{thm:frontier-single-flip}, now once in the primal block and
once in the dual block, gives
\begin{equation}
 \left\|
 \begin{bmatrix}
  \mathcal A_\varsigma(X_w)-b\\
  \operatorname{svec}(\mathcal A_\varsigma^*(y_w)+S_w-C)
 \end{bmatrix}\right\|_2
 \leq\frac{2B_{\rm KKT}H\sqrt{2s_{\rm KKT}M_{\rm dep}}}{N}.
                                                               \label{eq:structural-sdp-single-flip}
\end{equation}
Keeping the dual multipliers signed gives the stronger weighted estimate
\begin{equation}
 \left\|
 \begin{bmatrix}
  \mathcal A_\varsigma(X_w)-b\\
  \operatorname{svec}(\mathcal A_\varsigma^*(y_w)+S_w-C)
 \end{bmatrix}\right\|_2^2
 \leq4(s_r+s_c)B^2H^2\sum_i M_iw_i^2,
 \label{eq:structural-sdp-weighted-residual}
\end{equation}
where $M_i$ counts dependent coefficient positions assigned to bit $i$.
For uniform weights this is
$4(s_r+s_c)B^2H^2M_{\rm dep}/N^2$, which implies
\eqref{eq:structural-sdp-single-flip}.
The PSD variables need no coordinatewise sign split; convexity preserves the
cone, and the identity block on $S$ is not input dependent.
For multibit coefficients the degree factor of
\cref{prop:frontier-multibit} enters, as in
\cref{sec:frontier-mixture-supplements}.

\begin{proposition}[SDP mixture phase boundary]
\label{prop:structural-sdp-phase}
Suppose $0<m\leq M$, every neighboring center obeys the common sandwich
$mI\preceq X_i\preceq MI$, and $R=M/m$.  Assume the one-bit-local dependence
stated before \eqref{eq:structural-sdp-single-flip}, that all $N$ bits are
sensitive, and that the selected neighboring centers have a common
mixture-stable wrong
output.  This may be either a fixed affine
score with a uniform signed margin, or a fixed symmetric contraction $O$
such that
\begin{equation}
 f(\xi)\inner O{X^\xi}\geq\gamma\tr X^\xi               \label{eq:structural-sdp-povm-margin}
\end{equation}
for an output $f\in\{-1,+1\}$ and a positive margin $0<\gamma\leq1$.
The contraction condition means $-I\preceq O\preceq I$; the binary
measurement has effects $(I+O)/2$ and $(I-O)/2$ on $X/\tr X$.
If a base-instance contract demands the correct
output from every positive-definite triple with combined residual at most
$\varepsilon\geq0$, algebraic gap $r\mu$, and either common-$\mu$ or point-centered
operator width at most $\theta$, then soundness requires
\begin{equation}
 \boxed{
 \frac{8s_{\rm KKT}B_{\rm KKT}^2H^2M_{\rm dep}}{N^2}>\varepsilon^2
 \quad\text{or}\quad
 \frac{(R-1)^2}{4R}>\theta.}                             \label{eq:structural-sdp-phase}
\end{equation}
For a Frobenius neighborhood of width $\eta_F$, the second alternative is
\begin{equation}
 \sqrt r\,\frac{(R-1)^2}{4R}>\eta_F,                     \label{eq:structural-sdp-frobenius-phase}
\end{equation}
unless \eqref{eq:structural-sdp-variance-frobenius} is smaller.
The residual alternative can be sharpened to
$4(s_r+s_c)B^2H^2M_{\rm dep}/N^2>\varepsilon^2$.
Under point centering the algebraic gap remains $r\mu$, not $r\mu_w$.
\end{proposition}

\begin{proof}
Equation \eqref{eq:structural-sdp-single-flip} gives the first certificate,
\eqref{eq:structural-sdp-gap} gives the second, and
\cref{thm:structural-sdp-kantorovich} gives operator width at most
$K(R)-1=(R-1)^2/(4R)$ under either centering convention.  The affine score is
preserved by averaging.  Under \eqref{eq:structural-sdp-povm-margin}, trace
weighting gives
\[
 -f(\varsigma)\tr\left(O\frac{X_w}{\tr X_w}\right)\geq\gamma,
\]
so the normalized density matrix has constant bias for the wrong sign.  If
neither strict inequality in \eqref{eq:structural-sdp-phase} held, this wrong
output would satisfy the contract.  Finally,
$\norm{Z_w-I}_F\leq\sqrt r\norm{Z_w-I}_{\rm op}$ gives
\eqref{eq:structural-sdp-frobenius-phase}; the same projection argument used
after \eqref{eq:structural-sdp-variance-operator} covers point centering.
\end{proof}

Since
\begin{equation}
 K(R)-1=\sinh^2\!\left(\frac{\log R}{2}\right),
\end{equation}
under a common sandwich $0<m\leq M$ and $mI\preceq X_i\preceq MI$, with
$R=M/m\geq1$, the worst-case common-$\mu$ Frobenius guarantee of width
$\eta_F\geq0$ based only on this global ratio is equivalent to
\begin{equation}
 \log R\leq2\operatorname{arsinh}(\sqrt{\eta_F}\,r^{-1/4})
 =O(r^{-1/4}).                                             \label{eq:structural-sdp-high-dim-law}
\end{equation}
Indeed, the equal mixture $X_1=mI$, $X_2=MI$ attains
$Z_w=K(R)I$ in every dimension.  The same ratio criterion is sufficient
for point centering, but this example has zero point-centered width, so
it does not establish necessity for that convention.  A finite consequence
of the criterion is $\log R\leq2\sqrt{\eta_F}\,r^{-1/4}$.
Actual Frobenius variance can be small even when the global spectral
ratio is not.

The common sandwich is load-bearing.  Individual condition numbers cannot
replace its ratio: the scalars $X_1=1,X_2=9$ each have condition number one,
but their equal mixture has $\norm{Z_w-I}_{\rm op}=16/9$.

\begin{remark}[Seven scope limits]
\label{rem:structural-sdp-mixture-scope}
The common-basis bounds $mI\preceq X_i\preceq MI$ are strong.  Operator
control incurs a generic $\sqrt r$ loss in Frobenius norm.  The factors
$m^{-2}$ in the additive-variance bounds are real.  Every sparsity and height
parameter depends on the chosen Frobenius-isometric $\operatorname{svec}$.
Output stability is a separate hypothesis, and
\eqref{eq:structural-sdp-povm-margin} says nothing about normalized
$\operatorname{svec}$ amplitudes.  The mixture is existential, so the theorem
is not a query lower bound without a reduction.  Finally,
\eqref{eq:structural-sdp-gap} is an algebraic identity at an infeasible point.
These qualifications are part of the theorem's usable boundary.

The arithmetic--harmonic operator inequality and its Kantorovich reverse are
classical
\cite{structural-kubo1980-means,structural-ando1983-arithmetic,
structural-marshall1990-kantorovich,lin2013-on-an-operator-kantorovich-inequality}.
The nearest slogan-level collision is conic-IPM warm starting
\cite{structural-chen2025-conic-warmstart}, which moves one old solution to a
centered start for a perturbed problem.  The conjunction here is different:
cross-instance averaging, a sparse incidence budget, simultaneous primal and
dual residuals, and a wrong-output boundary.  The scalar LP corollary was
already proved in \cref{prop:frontier-central-dichotomy}; the present result is
the SDP operator theory behind it.
\end{remark}

\paragraph{Formal verification.}
The Lean development in \path{formal/QipmFormal/SDPMixture/} verifies the
finite matrix-mixture argument in this subsection.  It includes the
variance identity and equality case, the sharp noncommutative sandwich,
operator and Frobenius variance bounds, trace recentering, and the
algebraic-gap discrepancy.  A constructed Frobenius-isometric symmetric
vectorization connects actual SDP measurement maps and their adjoints to
weighted, uniform, and multibit residual bounds.  The soundness results
combine these properties with explicit affine or trace-normalized
observable output contracts under both centering conventions.  The
scalar ratio criterion, endpoint sharpness, and the distinction between
the two centering conventions are also checked.  The claim map,
qualifications, and reproduction command are in \path{formal/SDP_MIXTURE.md}.
The development permits only \texttt{propext}, \texttt{Classical.choice},
and \texttt{Quot.sound}.  It does not assert existence of the neighboring
centers for arbitrary inputs or formalize an IPM trajectory or quantum
query reduction.

\subsection{Rank required for SDP complementarity progress}
\label{sec:structural-rank-progress}

Low rank at an SDP optimum does not imply low rank of an accurate Newton
direction.  The clean statement is in Nesterov--Todd (NT) coordinates.  At
an exact center $XS=\mu I$, let $W$ be the NT scaling, so $WSW=X$, put
$P=W^{-1/2}$, and define
\begin{equation}
 U=P\Delta XP^T,\qquad V=P^{-T}\Delta SP^{-1}.
\end{equation}
Here $U$ denotes the NT-scaled primal direction and is unrelated to the
row-change matrix and spectral endpoint denoted by $U$ above.
Write the symmetric scaled complementarity equation operationally as
\begin{equation}
 \sqrt\mu(U+V)=-(1-\sigma)\mu I+\widetilde R.             \label{eq:structural-scaled-complementarity}
\end{equation}
This definition avoids dependence on competing residual conventions.

\begin{theorem}[Sharp rank--progress barrier]
\label{thm:structural-rank-progress}
Let $0\leq\sigma<1$.  For $1\leq p<\infty$, if
\begin{equation}
 \norm{\widetilde R}_{S_p}
 \leq\theta(1-\sigma)\mu n^{1/p},\qquad 0\leq\theta<1,
                                                               \label{eq:structural-rank-residual}
\end{equation}
then
\begin{equation}
 \boxed{\rank(\Delta X)+\rank(\Delta S)
 \geq\left\lceil(1-\theta^p)n\right\rceil.}             \label{eq:structural-rank-bound}
\end{equation}
For spectral error,
\begin{equation}
 \norm{\widetilde R}_{\mathrm{op}}<(1-\sigma)\mu
 \quad\Longrightarrow\quad
 \rank(\Delta X)+\rank(\Delta S)\geq n.                 \label{eq:structural-rank-spectral}
\end{equation}
The result is invariant under primal congruence and holds blockwise on a
product PSD cone.
\end{theorem}

\begin{proof}
Congruence preserves rank, so, with
$k=\rank(\Delta X)+\rank(\Delta S)$,
\[
 \rank(U+V)\leq\rank(U)+\rank(V)=k.
\]
Equation \eqref{eq:structural-scaled-complementarity} makes
$\widetilde R$ the error of approximating $(1-\sigma)\mu I$ by a
rank-at-most-$k$ matrix.  Eckart--Young--Mirsky gives
\[
 \norm{\widetilde R}_{S_p}
 \geq(1-\sigma)\mu(n-k)_+^{1/p}.
\]
Combining this with \eqref{eq:structural-rank-residual} proves
\eqref{eq:structural-rank-bound}.  If $k<n$, a kernel vector of $U+V$
sees residual exactly $(1-\sigma)\mu$, proving
\eqref{eq:structural-rank-spectral}.
\end{proof}

The bound is sharp as a matrix statement: take
$U=-(1-\sigma)\sqrt\mu\diag(I_k,0)$ and $V=0$.  The residual is precisely
$(1-\sigma)\mu\diag(0,I_{n-k})$.  The approximation theorem used in the
proof is classical \cite{structural-eckart1936-low-rank,
structural-mirsky1960-unitarily-invariant}; the rank--progress interpretation
is the structural content here.

\subsubsection{A path witness and the norm audit}

Let $L$ be the Laplacian of the $n$-vertex path and consider
\begin{equation}
 \min\ \inner L X\qquad\text{subject to}\qquad
 \tr X=1,\quad X\succeq0.                                \label{eq:structural-path-sdp}
\end{equation}
For $u=\mathbf1/\sqrt n$, the unique optimum is $X^*=uu^T$, while the
dual optimum has slack $S^*=L$; hence strict complementarity holds.  Its
exact central path is
\begin{equation}
 y(t)=-t,\quad S(t)=L+tI,\quad Z(t)=\tr(S(t)^{-1}),
 \qquad \mu(t)=Z(t)^{-1},\quad X(t)=\mu(t)S(t)^{-1}.      \label{eq:structural-path-center}
\end{equation}
With $T(t)=\tr(S(t)^{-2})$, differentiation gives
\begin{equation}
 \dot S=I,\qquad \dot\mu=\mu^2T,\qquad
 \dot X=\dot\mu S^{-1}-\mu S^{-2}.                       \label{eq:structural-path-tangent}
\end{equation}
If $0=\lambda_0<\lambda_1\leq\cdots$ are the eigenvalues of $L$, the
eigenvalues of $\dot X$ are
\begin{equation}
 \frac{\mu}{(t+\lambda_i)^2}
 \left[\frac TZ(t+\lambda_i)-1\right].                   \label{eq:structural-path-raw-spectrum}
\end{equation}
The zero-mode entry is negative.  Every other entry is positive whenever
$t\leq\lambda_1/(2\sqrt n)$, because
$T/Z>1/(t+\lambda_1)$ follows from
$\lambda_1^2+\lambda_1t-(n-1)t^2>0$.  Thus both tangent matrices are full
rank arbitrarily close to a rank-one optimum.  The exact Newton direction
targeting $\sigma\mu$ is the scalar multiple
\begin{equation}
 (\Delta X,\Delta S,\Delta y)
 =-\frac{(1-\sigma)\mu}{\dot\mu}(\dot X,\dot S,\dot y).
                                                               \label{eq:structural-path-newton}
\end{equation}

Full algebraic rank alone is weak, so the scaled spectrum matters.  Put
$d_i=t+\lambda_i$ and $\beta=T/Z$.  Since
$W=\sqrt\mu S^{-1}$ on \eqref{eq:structural-path-center}, the NT-scaled
tangent eigenvalues are
\begin{equation}
 u_i=\sqrt\mu\left(\beta-d_i^{-1}\right),\qquad
 v_i=\sqrt\mu d_i^{-1},\qquad u_i+v_i=\sqrt\mu\beta.
                                                               \label{eq:structural-path-nt-spectrum}
\end{equation}
If $t\leq\lambda_1/(4\sqrt n)$, then $\beta d_i\geq3$ for $i\geq1$:
\[
 \beta\lambda_1
 \geq\frac{\lambda_1/t}{1+(n-1)t/\lambda_1}>3.
\]
Hence
\begin{equation}
 \frac23\sqrt\mu\beta<u_i<\sqrt\mu\beta\qquad(i\geq1),
                                                               \label{eq:structural-path-flat}
\end{equation}
and every rank-$k<n-1$ approximation obeys
\begin{equation}
 \norm{U-U_k}_F\geq\frac23\sqrt\mu\beta\sqrt{n-1-k}.
                                                               \label{eq:structural-path-approx-rank}
\end{equation}
Using the path trace identities below to control the exceptional zero mode,
any fixed relative Frobenius target error below $2/3$ therefore requires
rank $\Omega(n)$ in the NT metric.  More generally, for any fixed target
error below one, shrinking the constant in
$t\leq\lambda_1/(4\sqrt n)$ makes all $n-1$ nonexceptional singular values
an arbitrarily large fixed fraction of $\sqrt\mu\beta$ and again forces
rank $\Omega(n)$.

The coordinate system is essential.  For the path,
\begin{equation}
 \lambda_k=4\sin^2\frac{\pi k}{2n},\quad
 \tr L^+=\frac{n^2-1}{6},\quad
 \tr(L^+)^2=\frac{2n^4+5n^2-7}{180}.                    \label{eq:structural-path-traces}
\end{equation}
As $t\downarrow0$, the raw tangent spectrum tends to
$\{-\tr L^+,\lambda_1^{-1},\ldots,\lambda_{n-1}^{-1}\}$, whose stable rank
tends to $7/5$ and whose best rank-one relative Frobenius error tends to
$\sqrt{2/7}$.  At $t=\lambda_1/(2\sqrt n)$, by contrast, the $n-1$
nonzero-mode NT singular values are asymptotically flat.  This regime is
late and ill conditioned:
\begin{equation}
 t=\frac{\pi^2}{2n^{5/2}}(1+o(1)),\qquad
 \cond(L+tI)=\frac8{\pi^2}n^{5/2}(1+o(1)).               \label{eq:structural-path-late}
\end{equation}
For smaller $t$ only the lower bound $\Omega(n^{5/2})$ is uniform.

\subsubsection{Why the barrier is self-limiting}
\label{sec:structural-rank-escapes}

The rank theorem obstructs low-rank \emph{factor tomography} of the
one-block NT direction.  It does not prove an $\Omega(n^2)$ storage bound,
and the path witness makes the escapes explicit.

First, chordal conversion replaces the single block by edge blocks
\begin{equation}
 Y_i=\begin{pmatrix}x_{ii}&x_{i,i+1}\\x_{i,i+1}&x_{i+1,i+1}\end{pmatrix}
 \succeq0\qquad(1\leq i<n),                              \label{eq:structural-path-blocks}
\end{equation}
with overlap equations $(Y_i)_{22}=(Y_{i+1})_{11}$ and trace equation
\begin{equation}
 (Y_1)_{11}+\sum_{i=1}^{n-1}(Y_i)_{22}=1.                \label{eq:structural-path-local-trace}
\end{equation}
The objective is the sum of inner products with
$\left[\begin{smallmatrix}1&-1\\-1&1\end{smallmatrix}\right]$.
Restriction of a global PSD matrix is feasible in this formulation.
Conversely, consistent PSD edge matrices have a PSD completion because the
path is chordal: construct Gram vectors successively, adding one orthogonal
coordinate per new vertex.  Objective and trace use only specified entries,
so the reformulation is exact.  Its barrier Hessian has constant-size blocks;
elimination gives a tridiagonal normal matrix with one dense border.  A
tridiagonal factorization and one scalar Schur complement use $O(n)$ memory
and arithmetic per Newton solve.  The conservative short-step total is
$O(n^{3/2}\log(n/\epsilon))$ arithmetic and $O(n)$ memory.  This is a
concrete instance of the formulation dependence emphasized in the chordal
conversion literature \cite{zhang2024-complexity-of-chordal-conversion-for}.

Second, the tangent itself has the implicit resolvent representation
\eqref{eq:structural-path-tangent}.  More generally, accurate low-rank SDP
methods may use low-rank endpoint structure to precondition or reformulate a
dense direction rather than factorizing that direction
\cite{bellavia2021-a-relaxed-interior-point-method,
zhang2024-well-conditioned-primal-dual-interior}.  These are compatible with
\cref{thm:structural-rank-progress}.

Third, high approximate matrix rank does not imply a hard amplitude state.
Writing $a=\sqrt\mu\beta$, the path estimates give
\begin{equation}
 \norm{U/a-(I-uu^T)}_{\rm op}=O(n^{-1/2}),\qquad
 \norm{U/a-(I-uu^T)}_F=O(1).                             \label{eq:structural-path-state-close}
\end{equation}
After normalization, $\operatorname{vec}(U)$ is $O(n^{-1/2})$-close to
\begin{equation}
 \frac{\ket I-\ket{uu^T}}{\sqrt{n-1}}.                  \label{eq:structural-path-lcu-state}
\end{equation}
For $u=\mathbf1/\sqrt n$, $\ket I/\sqrt n$ is a maximally entangled index
state and $\ket{uu^T}$ is a product of uniform states.  A two-term LCU
prepares \eqref{eq:structural-path-lcu-state} with constant success
probability and $\widetilde O(\log n)$ gates.  The theorem therefore makes no
state-preparation lower-bound claim.

Boundary analyticity under strict complementarity is classical
\cite{halicka2002-analyticity-of-the-central-path}.  The rank law is a sharp
local obstruction, but cone reformulation, implicit resolvents, and compact
state-preparation circuits are first-class escapes rather than footnotes.

\subsection{Sparse support and sparse-system embeddings}
\label{sec:structural-sparse-results}

The next statements concern standard-form complementarity, not arbitrary
optimization algorithms.  At an exact LP center $XSe=\mu e$, let
$0\leq\sigma<1$.  An inexact centering direction satisfies
\begin{equation}
 S\Delta x+X\Delta s=-(1-\sigma)\mu e+r,\qquad
 \norm r_2\leq\eta\mu.                                   \label{eq:structural-inexact-centering}
\end{equation}
Let $\operatorname{psupp}(\Delta x,\Delta s)$ be the indices where at least
one member of the coordinate pair is nonzero.

\begin{theorem}[Complementarity support barrier]
\label{thm:structural-support-barrier}
For $0\leq\sigma<1$, every direction satisfying
\eqref{eq:structural-inexact-centering} obeys
\begin{equation}
 \boxed{|\operatorname{psupp}(\Delta x,\Delta s)|
 \geq n-\frac{\eta^2}{(1-\sigma)^2}.}                    \label{eq:structural-support-count}
\end{equation}
For a short step $1-\sigma=a/\sqrt n$ with $\eta<a$, this is
$[1-(\eta/a)^2]n$.  If a $k$-sparse latent vector $z$ is reconstructed as
$(\Delta x,\Delta s)=Bz$, and each column of $B$ touches at most $g$
complementarity pairs, then
\begin{equation}
 \boxed{gk\geq n-\frac{\eta^2}{(1-\sigma)^2}.}            \label{eq:structural-locality-product}
\end{equation}
\end{theorem}

\begin{proof}
Let $Z$ contain the indices at which both update coordinates vanish.  For
each $i\in Z$, \eqref{eq:structural-inexact-centering} forces
$r_i=(1-\sigma)\mu$.  Hence
\[
 |Z|(1-\sigma)^2\mu^2\leq\norm r_2^2\leq\eta^2\mu^2,
\]
which proves \eqref{eq:structural-support-count}.  The support of $Bz$ is
contained in the union of at most $k$ column pair-supports, of total size at
most $gk$, proving \eqref{eq:structural-locality-product}.
\end{proof}

The count is realized by a sparse, incompressible direction.  For $n\geq3$,
consider
\begin{equation}
 \min x_2\quad\text{subject to}\quad
 x_1+x_2=1,\quad x_{i-1}-x_i=0\ (3\leq i\leq n),\quad x\geq0.
                                                               \label{eq:structural-incompressible-lp}
\end{equation}
Its equality matrix has maximum row and column sparsity two (the endpoint
columns have only one nonzero), and its unique
optimum is $e_1$.  The feasible segment is
$x(t)=(1-t,t,\ldots,t)$, and, for $0<t<(n-1)/n$, the central parameter is
\begin{equation}
 \mu(t)=\frac{t(1-t)}{n-1-nt}.
\end{equation}
Indeed, eliminating the equalities gives the one-variable barrier
$t-\mu[\log(1-t)+(n-1)\log t]$; its stationarity equation is the displayed
formula.
At $t=1/2$, every nonzero central tangent is proportional to
$(-1,1,\ldots,1)$.  Its normalized best $k$-sparse approximation error is
exactly $\sqrt{1-k/n}$; a sparse optimum therefore does not make the central
direction compressible.

There is also no hidden geometric restriction on a positive direction
profile.  Given any $u\in\R_{++}^n$ with $\norm u_2=1$, choose $\mu_0>0$ so
that $u_i^{-2}>4\mu_0$ and set
\begin{equation}
 c_1=-\sqrt{u_1^{-2}-4\mu_0},\qquad
 c_i=\sqrt{u_i^{-2}-4\mu_0}\quad(i\geq2).                \label{eq:structural-universal-cost}
\end{equation}
For the bound-constrained QP
\begin{equation}
 \min_{x\geq0}\ \tfrac12\norm x_2^2+c^Tx,               \label{eq:structural-universal-qp}
\end{equation}
the unique optimum is the one-sparse vector $(-c_1)e_1$, with strict
complementarity.  Central stationarity gives $s=x+c$ and
\begin{equation}
 x_i(\mu)=\frac{-c_i+\sqrt{c_i^2+4\mu}}2,\qquad
 x_i'(\mu)=\frac1{\sqrt{c_i^2+4\mu}},
\end{equation}
so $x'(\mu_0)=u$.  The row-scaled complementarity Jacobian is
$\diag(2x(\mu_0)+c)=\diag(1/u_i)$, which is one-sparse and constant
conditioned for near-uniform $u$.  Combining this realization with coherent
tomography gives the established
$\widetilde\Omega(n/\epsilon)$ solution-unitary readout bound on the
near-uniform hard family \cite{joran2023-quantum-tomography-state-preparation-unitaries}.
It is a modular readout bound: an algorithm allowed to inspect the explicit
coefficients \eqref{eq:structural-universal-cost} can learn $u$ there.

The condition-one box-LP family behind the older version of this observation
is superseded in strength by the original-primal and bounded-data families
of \cref{sec:lp-prefix-rigidified,sec:lp-bounded-data}.  Its remaining lesson
is only that perfect reduced conditioning does not remove coefficient-query
or classical-output cost.  Likewise, the older coordinate-reset ledger is
superseded by the projective path theorem of \cref{sec:positive-lazy}; the
QRAM-write limitation is stated once in
\cref{sec:structural-geometric-locality}.

\subsubsection{A query-preserving OSS progress step}

The support theorem says what a physical update must touch.  A distinct
question is whether standard sparse-QLS lower bounds occur inside a valid
QIPM Newton solve.  They do.  Let $H\in\R^{n\times n}$ be real symmetric,
invertible, with $\norm H_2=1$ and bidirectional row/column-slot sparse
access.  Let $d$ be the declared slot bound; declared candidate slots may
contain zeros.  Let $\norm q_2=1$.  Fix $\theta\in(0,1)$, choose
\begin{equation}
 0<\delta\leq\min\{\theta,1/2\},\quad
 0\leq\alpha\leq\delta/4,\quad
 a=\delta q/\sqrt2,\quad D=\diag(a),\quad
 \sigma=1-\alpha/\sqrt n.                                \label{eq:structural-oss-parameters}
\end{equation}
Consider
\begin{equation}
 \min c^Tx\quad\text{subject to}\quad
 Ax=0,\ x\geq0,\qquad
 A=\begin{bmatrix}H&-H\end{bmatrix},\quad c=(e+a,e-a).    \label{eq:structural-oss-lp}
\end{equation}
At $x^0=e_{2n},y^0=0,s^0=c$, the complementarity average is one and the
$N_2$ residual norm is $\delta$, while the $N_\infty$ residual is at most
$\delta/\sqrt2$.  Thus the point belongs to both radius-$\theta$
neighborhoods.

\begin{theorem}[Feasible OSS embedding]
\label{thm:structural-oss-embedding}
With null-space basis $V=(I,I)^T$ and $r=1-\sigma$, the feasible OSS step is
\begin{equation}
 \underbrace{\begin{bmatrix}-H&I+D\\H&I-D\end{bmatrix}}_{M_{\rm OSS}}
 \binom{\Delta y}{\lambda}
 =\binom{-re-a}{-re+a}.                                  \label{eq:structural-oss-system}
\end{equation}
It has row sparsity at most $d+1$, column sparsity at most $2d$, and
$\cond(M_{\rm OSS})=\Theta_\delta(\cond(H))$.  Its exact solution is
\begin{equation}
 \lambda=-re,\qquad \Delta y=\sigma H^{-1}a.             \label{eq:structural-oss-solution}
\end{equation}
The normalized solution has probability at least $2/3$ in the $\Delta y$
block.  The physical step lands at
\begin{equation}
 x^1=\sigma e,\qquad s^1=(e+ra,e-ra),\qquad \mu^1=\sigma,
                                                               \label{eq:structural-oss-endpoint}
\end{equation}
with centrality-residual norm $\sigma r\delta$.  Every sparse query to the
LP or OSS data is simulated with $O(1)$ queries to $H$ and evaluations of the
known $q$.
\end{theorem}

\begin{proof}
Primal and dual feasibility at the start are immediate, and
$\norm{X^0s^0-e}_2=\norm{(a,-a)}_2=\delta$.  Substitution of
$\Delta x=V\lambda$ and $\Delta s=-A^T\Delta y$ into the centering equation
gives \eqref{eq:structural-oss-system}.  The orthogonal row transform
\[
 Q=2^{-1/2}\begin{bmatrix}-I&I\\I&I\end{bmatrix}
\]
gives
\begin{equation}
 QM_{\rm OSS}=\sqrt2\begin{bmatrix}H&-D\\0&I\end{bmatrix}.
\end{equation}
The triangular factor has constant norm and inverse
$\left[\begin{smallmatrix}H^{-1}&H^{-1}D\\0&I\end{smallmatrix}\right]$;
since $\norm D\leq\delta/\sqrt2$, its inverse norm is
$\Theta_\delta(\norm{H^{-1}})$.

Applying $Q$ to the right-hand side gives
$H\Delta y-D\lambda=a$ and $\lambda=-re$.  Since $De=a$, this proves
\eqref{eq:structural-oss-solution}.  Also
$\norm{\Delta y}_2^2\geq\sigma^2\delta^2/2$, whereas
$\norm\lambda_2^2=\alpha^2$, so the block probability is at least
\[
 \frac{\sigma^2\delta^2/2}{\sigma^2\delta^2/2+\alpha^2}>\frac23.
\]
Direct substitution proves \eqref{eq:structural-oss-endpoint} and gives
$X^1s^1-\mu^1e=\sigma r(a,-a)$.  All remaining entries and sparse positions
are fixed, proving the oracle simulation.
\end{proof}

After symmetric dilation, the theorem transfers the
$\Omega(\kappa\sqrt d)$ fixed-error and
$\Omega(\kappa\log(1/\epsilon))$ constant-sparsity QLS lower bounds of
Mori et al.\ to a strictly feasible short progress step
\cite{mori2026-sparsity-dependent-complexity-lower-bound}.  It is an
embedding of known QLS restrictions, not a new QLS technique or an
end-to-end LP lower bound.  Indeed, the physical $(\Delta x,\Delta s)$ above
is explicit; the hard block is $\Delta y$, and algorithms that bypass
accurate OSS-state preparation are outside the statement.
Here $d$ is the declared sparse-oracle slot bound, which may count padded
zero slots.  A symmetric dilation of a nonsymmetric source matrix requires
bidirectional row- and column-slot access.  Mori et al.'s clock family has
known successor/predecessor slot locations and reversible hidden updates, so
both directions are simulated with constant query overhead.

\subsubsection{What bounded treewidth preserves}

The same precision restriction survives at treewidth two.  Mori et al.'s
hard matrix has the form
\begin{equation}
 A=I-\rho B,                                               \label{eq:structural-clock-matrix}
\end{equation}
where $B$ is a clock permutation and every row and column has one identity
and one permutation position.  Form its symmetric dilation
\begin{equation}
 K_A=\begin{bmatrix}0&A^T\\A&0\end{bmatrix},\qquad
 \widehat b=\binom0b.                                    \label{eq:structural-treewidth-dilation}
\end{equation}

\begin{theorem}[Treewidth-two KKT precision]
\label{thm:structural-treewidth-two}
There are real symmetric, indefinite, two-sparse equality-KKT matrices $K_A$
whose support graphs have treewidth at most two and for which preparing an
$\epsilon$-approximation to the normalized solution requires
\begin{equation}
 \Omega\!\left(\cond(K_A)\log(1/\epsilon)\right)          \label{eq:structural-treewidth-qls}
\end{equation}
queries in the parameter range of the underlying QLS theorem.
\end{theorem}

\begin{proof}
One has
$K_A^{-1}\widehat b=(A^{-1}b,0)^T$, and the singular values of $K_A$ are
those of $A$, each repeated.  The identity positions of $A$ give one perfect
matching in the bipartite support graph of $K_A$, while the clock-permutation
positions give a second.  Their union is a disjoint union of paths and even
cycles, so its treewidth is at most two.  Successor and predecessor clock
queries simulate every sparse query with constant overhead.  The state and
condition number are preserved, so the QLS precision lower bound transfers.
Finally, $K_A$ is the equality-KKT matrix of $\min0$ subject to $Ax=b$.
\end{proof}

The two-sparse clock dilation of Wang--Zhang has the same low-treewidth
support phenomenon but proves a different query-depth statement and not the
$\log(1/\epsilon)$ precision factor
\cite{wang2024-tight-quantum-depth-lower-bound}.
This does not say that bounded-treewidth optimization is generically hard.
The stronger classical fact points the other way: supplied width-$\tau$
elimination gives $O(N\tau^2)$ scalar Cholesky work and $O(N\tau)$ storage
for an $N$-dimensional SPD Newton system.  A dense classical output already
costs $\Omega(N)$, so the dimension-only full-output speedup ceiling is
$O(\tau^2)$.  Nearly linear bounded-treewidth IPMs sharpen this baseline
\cite{dong2020-a-nearly-linear-time-algorithm,gu2022-a-faster-small-treewidth-sdp}.

There is nevertheless scalar LP value hardness at an exact chosen width.
For Boolean matrices $Z^i\in\{0,1\}^{c\times b}$, start from the pre-dual
program
\begin{equation}
 \begin{aligned}
  \max_{w,v^1,\ldots,v^a}\quad&\sum_{k=1}^c w_k\\
  \text{subject to}\quad&
  w-\sum_{i=1}^a Z^iv^i\leq0,\qquad
  e^Tv^i\leq1\ (1\leq i\leq a),\\
  &w,v^i\geq0.
 \end{aligned}                                            \label{eq:structural-apers-predual}
\end{equation}
Its value is
\begin{equation}
 \sum_{i=1}^a\max_{j\in[b]}\sum_{k=1}^c Z_{ijk}.          \label{eq:structural-apers-value}
\end{equation}
Add a nonnegative, zero-objective $z_i$ to every top inequality and bottom
inequality $i$, always with coefficient $+1$.

\begin{theorem}[LP value hardness at exact treewidth]
\label{thm:structural-exact-treewidth}
For integers $a,b,c\geq1$, the augmented LP has
$p=c+a(b+1)$ variables, $q=c+a$ constraints, and
constraint-intersection treewidth exactly $\tau=c$.  When $c\leq a$,
constant-additive approximation of its scalar value requires
\begin{equation}
 \Omega(a\sqrt b\,c)=\Omega(\tau\sqrt{pq})                \label{eq:structural-exact-treewidth-value}
\end{equation}
fixed-position quantum coefficient queries.  The corresponding randomized
bound is $\Omega(abc)=\Omega(p\tau)$.
\end{theorem}

\begin{proof}
For fixed $v^i$, each top constraint is tight at an optimum, so summing the
coordinates of $w$ separates the $a$ simplex problems in $v^i$ and gives
\eqref{eq:structural-apers-value}.  Setting $z=0$ extends every original
feasible point.  Conversely, deleting positive $z$ from an augmented point
only relaxes its inequalities, so the value is unchanged.

Let $C$ be the $c$ top-constraint vertices.  The column of $z_i$ makes
$C\cup\{i\}$ a clique, and no variable column contains two bottom vertices.
Thus the intersection graph is exactly $K_c\vee I_a$, with bags
$C\cup\{i\}$, and has treewidth $c$.  The dimensions follow by counting
$w,v,z$.  The published Boolean reduction for
\eqref{eq:structural-apers-value} gives
$\Omega(a\sqrt b\,c)$ quantum and $\Omega(abc)$ randomized fixed-position
queries.  Since
\[
 a\sqrt b\,c=\tau\sqrt{(q-\tau)(p-q)},\qquad
 abc=\tau(p-q),
\]
and $c\leq a$ makes $q-\tau=\Theta(q)$ and $p-q=\Theta(p)$, the displayed
bounds follow.
\end{proof}

This is a reparameterized structural corollary of an existing reduction
\cite{apers2026-quantum-speedups-linear-programming-interior}, not a new
adversary argument.  It applies to individual coefficient queries; a whole
column reveals $\Theta(\tau)$ coefficients, and a whole top row can reveal
far more.  The broad claim that treewidth was absent from QLS lower-bound
work would also be false in view of Khanpour--Talkington
\cite{khanpour2026-proving-limits-power-flow}.  The new restriction here is
only the treewidth-two KKT precision embedding.

\subsection{Geometric locality of central-point access}
\label{sec:structural-geometric-locality}

Sparse-query access suppresses where the input is stored.  On bounded-range
hardware, that geometry can itself be an obstruction.  Let hidden signs
$\varsigma_1,\ldots,\varsigma_N\in\{\pm1\}$ and introduce
$d_i=u_i-v_i$ with $u_i,v_i\geq0$.  For fixed $\alpha\in(0,1)$, consider
\begin{equation}
 \begin{aligned}
  \min\quad&\sum_{i=0}^N(u_i+v_i)+\alpha(u_N-v_N),\\
  \text{subject to}\quad&d_0=1,\\
  &d_i-\varsigma_i d_{i-1}=0\quad(1\leq i\leq N).
 \end{aligned}                                             \label{eq:structural-locality-lp}
\end{equation}
Every row has at most four nonzeros, every column at most two, and the
constraint-intersection graph is a path.  Put
$p_i=\prod_{j=1}^i\varsigma_j$ and $q_i=u_i+v_i$.  Feasibility gives
$d_i=p_i$.  Taking every $q_i=2$ is strictly primal feasible; with the
convention $A^Ty+s=c$, the choice $y=0$ gives $s=c>0$ and strict dual
feasibility.  The barrier restricted to the feasible affine space is
\begin{equation}
 \sum_{i=0}^Nq_i+\alpha p_N
 -\mu\sum_{i=0}^N\log\frac{q_i^2-1}{4}.                  \label{eq:structural-locality-barrier}
\end{equation}
Thus every central coordinate equals
\begin{equation}
 q(\mu)=\mu+\sqrt{\mu^2+1}.                              \label{eq:structural-locality-center}
\end{equation}
At $\mu=3/4$, $q=2$ and
$(u_i,v_i)=(3/2,1/2)$ or $(1/2,3/2)$ according to $p_i$.  Reading the final
cell with coordinate error below $1/2$ therefore reveals parity.

This information cost is invisible to reduced conditioning.  The public
vectors $V_i=(e_{u_i}+e_{v_i})/\sqrt2$ form an orthonormal null-space basis,
and for the full logarithmic-barrier Hessian $H_x=X^{-2}$,
\begin{equation}
 V^TH_xV
 =2\bigl[(q+1)^{-2}+(q-1)^{-2}\bigr]I.                   \label{eq:structural-locality-condition-one}
\end{equation}
It has condition number one for every $\mu>0$.

Place input register $\varsigma_i$ at a site $v_i$ of a hardware graph
$\Gamma$.  A depth-$T$, range-$r$ circuit has gates of hardware diameter at
most $r$ in each layer.  Its fixed initial workspace may contain arbitrary
input-independent entanglement.  A measurement at site $o$ is locally
parity-readable if it returns $p_N$ with probability at least $2/3$ for every
input.

\begin{theorem}[Central-oracle light cone]
\label{thm:structural-light-cone}
Let $R_\Gamma(o;v)=\max_{1\leq i\leq N}\operatorname{dist}_\Gamma(o,v_i)$
and let $b=\max_{w\in\Gamma}|\{i:v_i=w\}|$ be the maximum site occupancy.
Every such circuit constructing a locally readable final central cell obeys
\begin{equation}
 \boxed{T\geq\left\lceil\frac{R_\Gamma(o;v)}{r}\right\rceil.}
                                                          \label{eq:structural-light-cone}
\end{equation}
Consequently, when at most $b$ input registers share a site, a path
requires $T\geq(\lceil N/b\rceil-1)/(2r)$ for every choice of output site,
which is $\Omega(N/(br))$ once $N\ge2b$, and a $d$-dimensional
nearest-neighbor grid requires $(2rT+1)^d\geq\lceil N/b\rceil$, hence
$T\ge((N/b)^{1/d}-1)/(2r)$, which is $\Omega((N/b)^{1/d}/r)$ once
$N\ge2^db$.
Pairwise-distinct placement is the case $b=1$.
\end{theorem}

\begin{proof}
Pull a local output observable backwards through the circuit.  Its support
can grow by hardware distance at most $r$ per layer, so after $T$ layers it
is contained in $B_\Gamma(o,rT)$.  If some $v_j$ lies outside this ball, two
inputs differing only at $j$ have identical reduced input states on the
backward light cone.  Their workspace state is also identical, including
all of its input-independent entanglement.  The local output distributions
are therefore identical, although flipping $\varsigma_j$ flips $p_N$; one
distribution cannot return both answers with probability $2/3$.

Thus the ball $B_\Gamma(o,rT)$ contains every occupied site, proving
\eqref{eq:structural-light-cone}.  Under occupancy at most $b$, at least
$\lceil N/b\rceil$ distinct sites are occupied.  On a path a radius-$rT$
ball has at most $2rT+1$ sites, and on the grid at most $(2rT+1)^d$, which
proves the two consequences.
\end{proof}

This is the standard backward-light-cone argument
\cite{barak2022-classical-algorithms-and-quantum-limitations}; the claimed
contribution is only its conjunction with the strictly feasible,
condition-one central-point witness.  It is not a lower bound in the
all-to-all or free-QRAM query model, for a homogeneous self-dual embedding
with a universal start, for an answer left in a global observable whose
measurement is declared free, or when unboundedly many input registers may
share one site.  The same lower bound holds classically, so it
is an obstruction to a claimed quantum depth advantage, not a
quantum--classical separation.  Input-dependent pre-entanglement can already
contain the answer and merely moves its construction cost into advice.

A stronger per-iteration QRAM-write claim does not follow.  A lazy oracle may
store the input once and evaluate a changed weight from the queried index and
a scalar path parameter.  An $\Omega(N)$ refresh theorem would need a
materialized-array or cell-probe model, or a lower bound for the composed
evaluation oracle.  Numerical change in all coordinates is not enough.

\subsection{Synthesis}
\label{sec:structural-synthesis}

These tools charge different resources and should not be collapsed into one
informal claim that sparsity fails.  Symmetry charges invariant virtual work
even after equality rows are mixed.  Central mixtures trade sparse
one-bit-local input incidence against operator spread.  Complementarity
progress charges pair
support or NT aggregate rank.  Sparse-QLS hardness survives a legitimate OSS
step and even a treewidth-two KKT representation, while exact treewidth also
admits strong classical algorithms.  Physical locality charges communication
only when the input and output are locally materialized with bounded per-site
occupancy.

Their escape routes are equally structural.  A nonlinear output may not be
mixture stable; central matrices may have a large common-basis spectral
spread; a nonlocal reconstruction map may keep a latent state sparse; chordal
conversion or an implicit resolvent may avoid matrix-factor tomography; an
algorithm may bypass the OSS; and a global or lazy oracle may avoid local
materialization.  The results above identify where such an escape must enter.
They do not supply a universal quantum query or runtime lower bound.
